\documentclass[12pt]{article}

\usepackage{amsmath}

\begin{document}

\section{Introduction}
\label{sec:introduction}

Planning is a core problem in artificial intelligence: given a model of an environment and a desired objective, an agent must construct a sequence of actions that leads from an initial state to a goal state while satisfying constraints and optimising some notion of cost or reward. Classical approaches formulate planning as search over a state space (e.g., A*), whereas reinforcement learning (RL) frames it as maximising cumulative reward through value functions or policies. More recently, a generative modelling perspective has emerged, in which trajectories themselves are treated as structured sequences that can be modelled probabilistically. From this viewpoint, planning becomes a problem of inference in a learned distribution over trajectories rather than explicit search in a hand-designed state graph.

In this project, planning in a 2D gridworld is formulated as autoregressive sequence modelling. A trajectory is represented as an ordered sequence of states (or state--action pairs), and a model is trained using next-step prediction on shortest paths generated by A*. Once trained, the model defines a learned probability distribution over feasible trajectories. Planning then corresponds to decoding: generating a trajectory from this distribution that satisfies the task objective.

The central question of this project is how different decoding strategies affect the quality, reliability, diversity, and computational cost of the generated plans. While the underlying model is fixed after training, inference procedures such as greedy decoding, stochastic sampling with temperature control, and beam search may extract substantially different trajectories from the same distribution. Greedy decoding prioritises local likelihood, potentially sacrificing global optimality or diversity. Stochastic sampling can increase diversity but may introduce instability or suboptimal paths. Beam search balances exploration and exploitation at the cost of increased computation. By systematically comparing these strategies, the project investigates how inference choices shape planning outcomes when planning is cast as generative sequence modelling.
\end{document}